\section{Apéndices}
\label{sec:apendice}

\subsection{Apéndice A. Glosario de Categorías Emocionales}

El esquema analítico multietiqueta implementado en SentiDataUMSA evalúa dieciocho categorías afectivas adaptadas a la interacción universitaria (Tabla~\ref{tab:glosario}).

\begin{table}[htbp]
  \centering
  \small
  \caption{Familias de Categorías Emocionales y Definición Operacional}
  \label{tab:glosario}
  \begin{tabular}{lp{9.5cm}}
    \toprule
    Familia & Categorías y Términos Representativos \\
    \midrule
    Positivas Sociales & Gratitud (\emph{gracias, bendiciones, auxi}), alegría (\emph{aprobé, genial}), satisfacción, confianza, entusiasmo y empatía. \\
    Negativas Académicas& Frustración (\emph{falla el sistema, no cargó}), ansiedad (\emph{plazos, nota}), enojo (\emph{estafa, injusto}), tristeza y decepción. \\
    Cognitivas / Neutras & Duda (\emph{alguien sabe, qué paralelo}), sorpresa, curiosidad y anticipación. \\
    \bottomrule
  \end{tabular}
  \vspace{4pt}
  \begin{minipage}{0.92\textwidth}
    \footnotesize
    \textit{Nota.} La asignación es probabilística multietiqueta y permite la coexistencia de estados afectivos mixtos en un mismo mensaje.
  \end{minipage}
\end{table}

\subsection{Apéndice B. Parámetros Globales de Ejecución}

La Tabla~\ref{tab:params} resume los hiperparámetros del sistema persistidos en la base de datos y en los archivos de configuración centralizados.

\begin{table}[htbp]
  \centering
  \small
  \caption{Parámetros de Configuración del Pipeline y Motores Analíticos}
  \label{tab:params}
  \begin{tabular}{lll}
    \toprule
    Parámetro Técnico & Valor Asignado & Justificación y Descripción \\
    \midrule
    \texttt{negation.window}        & 3          & Ventana sintáctica previa de tokens evaluados \\
    \texttt{sentiment.min\_evidence}& 1          & Umbral $\tau$ de evidencia para descartar indeterminación \\
    \texttt{sentiment.emoji\_weight}& 1,0        & Ponderación $\alpha$ de valencia en emojis del texto \\
    \texttt{sentiment.reaction\_wt} & 0,0        & Ponderación desacoplada para reacciones interactivas \\
    \texttt{topics.k}               & 5          & Número de tópicos latentes evaluados por Gibbs Sampling \\
    \texttt{markov.order}           & 1          & Orden estocástico de la cadena de estados polares \\
    \texttt{db.pool\_size}          & 5          & Conexiones concurrentes recomendadas en Neon PostgreSQL \\
    \bottomrule
  \end{tabular}
\end{table}

\subsection{Apéndice C. Modelo de Transición Estocástica (TikZ)}

La Figura~\ref{fig:grafo_markov_tikz} ilustra la máquina de estados estocástica finita de Cadenas de Markov estimada por máxima verosimilitud a partir de las transiciones entre mensajes consecutivos del corpus.

\begin{figure}[htbp]
  \centering
  % =============================================================================
% grafo_markov.tex — Diagrama Estocástico de Transición de Estados (TikZ)
% =============================================================================
\shorthandoff{><}
\begin{tikzpicture}[
    arrows={-Stealth},
    node distance=3.2cm,
    state/.style={circle, draw=blue!70!black, fill=blue!10, thick, minimum size=1.8cm, font=\bfseries\small, align=center},
    state_pos/.style={circle, draw=green!60!black, fill=green!15, thick, minimum size=1.8cm, font=\bfseries\small, align=center},
    state_neg/.style={circle, draw=red!70!black, fill=red!15, thick, minimum size=1.8cm, font=\bfseries\small, align=center},
    state_neu/.style={circle, draw=gray!70!black, fill=gray!15, thick, minimum size=1.8cm, font=\bfseries\small, align=center},
    prob/.style={font=\footnotesize\itshape, fill=white, inner sep=1.5pt, rounded corners=2pt}
]

    % Nodos de estados
    \node[state_pos] (POS) at (0, 2.5) {Positivo \\ \scriptsize ($\pi=11,5\%$)};
    \node[state_neg] (NEG) at (6, 2.5) {Negativo \\ \scriptsize ($\pi=5,4\%$)};
    \node[state_neu] (NEU) at (3, -1.8) {Neutro \\ \scriptsize ($\pi=83,1\%$)};

    % Bucles autorregresivos
    \path[thick] (POS) edge [loop left] node [prob] {$p = 0,160$} (POS);
    \path[thick] (NEG) edge [loop right] node [prob] {$p = 0,075$} (NEG);
    \path[thick] (NEU) edge [loop below] node [prob] {$p = 0,839$} (NEU);

    % Transiciones entre estados
    \path[thick, bend left=15] (POS) edge node [prob, above] {$0,064$} (NEG);
    \path[thick, bend left=15] (NEG) edge node [prob, below] {$0,102$} (POS);

    \path[thick, bend left=20] (POS) edge node [prob, left] {$0,776$} (NEU);
    \path[thick, bend left=20] (NEU) edge node [prob, right] {$0,109$} (POS);

    \path[thick, bend left=20] (NEG) edge node [prob, right] {$0,823$} (NEU);
    \path[thick, bend left=20] (NEU) edge node [prob, left] {$0,051$} (NEG);

\end{tikzpicture}
\shorthandon{><}

  \vspace{4pt}
  \caption{Grafo de Transición de Estados Emocionales y Probabilidades Estacionarias}
  \label{fig:grafo_markov_tikz}
\end{figure}

\subsection{Apéndice D. Menciones Docentes y Resolución Semestral}

La Tabla~\ref{tab:menciones_docentes} resume los resultados del módulo de resolución de entidades docentes y materias contra la oferta académica semestral UMSA.

\begin{table}[htbp]
  \centering
  \small
  \caption{Ranking de Docentes Más Mencionados, Materias Asociadas y Polaridad Observada}
  \label{tab:menciones_docentes}
  \resizebox{\textwidth}{!}{
  \begin{tabular}{llrrrl}
    \toprule
    Identificador & Nombre Docente & Menciones & \% Positivo & \% Negativo & Materia Principal \\
    \midrule
    DOC\_01 & Lic. Marcelo Balboa     & 342 & 68,4 & 14,2 & INF-131: Algoritmos \\
    DOC\_02 & M.Sc. Carmen Velasco    & 289 & 74,1 & 8,6  & INF-141: Base de Datos \\
    DOC\_03 & Lic. Rolando Huanca     & 245 & 55,2 & 28,1 & INF-111: Introducción \\
    DOC\_04 & Ing. Edgar Clavijo      & 218 & 62,0 & 16,5 & INF-121: Programación \\
    DOC\_05 & Lic. Patricia Condori   & 194 & 71,8 & 9,3  & INF-161: Minería de Datos \\
    DOC\_06 & Ing. Ramiro Quispe      & 165 & 49,4 & 32,1 & INF-144: Sistemas Operativos \\
    DOC\_07 & M.Sc. Fernando Terán    & 142 & 66,2 & 12,0 & INF-153: Redes \\
    \bottomrule
  \end{tabular}
  }
  \vspace{4pt}
  \begin{minipage}{0.94\textwidth}
    \footnotesize
    \textit{Nota.} Resolución automática ejecutada por \texttt{detectar\_menciones\_umsa.py} cruzando expresiones regulares compiladas (\texttt{alias\_regex}) contra la tabla \texttt{oferta\_academica\_semestral} (SCD Tipo 2). La tasa residual corresponde a menciones catalogadas como neutras o indeterminadas.
  \end{minipage}
\end{table}


\subsection{Apéndice E. Ficha Técnica General del Corpus}

La Tabla~\ref{tab:resumen_general} recapitula la ficha técnica consolidada del corpus y la infraestructura desarrollada.

\begin{table}[htbp]
  \centering
  \small
  \caption{Ficha Técnica Resumida del Proyecto ASDAT-253 y del Corpus UMSA}
  \label{tab:resumen_general}
  \resizebox{\textwidth}{!}{
  \begin{tabular}{ll}
    \toprule
    Parámetro de Investigación & Valor / Especificación Técnica \\
    \midrule
    Población / Entorno de estudio      & Grupo oficial Telegram de gestión académica (UMSA Informática) \\
    Ventana temporal analizada          & 15 de mayo de 2023 al 01 de marzo de 2026 (34 meses) \\
    Volumen de mensajes analizados      & 23.371 mensajes históricos \\
    Usuarios participantes identificados& 2.051 remitentes únicos \\
    Motor de base de datos              & Neon Serverless PostgreSQL 17/18 con SSL y pooling PgBouncer \\
    Tablas relacionales en esquema      & 16 tablas con soporte para CDC y JSONB \\
    Categorías emocionales evaluadas    & 18 emociones discretas con método léxico multietiqueta \\
    Términos en matriz TF-IDF           & 122.971 registros ponderados \\
    Modismos locales incorporados       & 538 jergas estudiantiles paceñas \\
    Número de tópicos latentes (LDA)    & $K = 5$ tópicos evaluados por muestreo de Gibbs \\
    Modelo estocástico de transiciones  & Cadena de Markov de 1.er orden sobre 3 estados polares \\
    Frameworks de software utilizados   & R 4.5 (\texttt{quanteda}, \texttt{topicmodels}) y Python 3.12 (Flask, Pandas, ECharts) \\
    \bottomrule
  \end{tabular}
  }
  \vspace{4pt}
  \begin{minipage}{0.94\textwidth}
    \footnotesize
    \textit{Nota.} Ficha consolidada generada automáticamente para anexos y memorias técnicas del curso DAT-253.
  \end{minipage}
\end{table}

